\chapter{Plano de validação}\label{plano-de-validauxe7uxe3o-e-crituxe9rios-de-sucesso}

{
\begingroup
\small
\singlespacing
\setlength{\tabcolsep}{4pt}
\begin{longtable}[]{@{}
  >{\raggedright\arraybackslash}p{(\linewidth - 4\tabcolsep) * \real{0.19}}
  >{\raggedright\arraybackslash}p{(\linewidth - 4\tabcolsep) * \real{0.385}}
  >{\raggedright\arraybackslash}p{(\linewidth - 4\tabcolsep) * \real{0.425}}@{}}
\toprule\noalign{}
\begin{minipage}[b]{\linewidth}\raggedright
Nível
\end{minipage} & \begin{minipage}[b]{\linewidth}\raggedright
Verificação
\end{minipage} & \begin{minipage}[b]{\linewidth}\raggedright
Critério de planejamento
\end{minipage} \\ \addlinespace[4pt]
\midrule\noalign{}
\endhead
\bottomrule\noalign{}
\endlastfoot
Matemático & Simetrias de Riemann, identidade de Bianchi, vínculos e
relação escalar & Igualdades simbólicas ou aritmética exata em
benchmarks. \\ \addlinespace[4pt]
Físico & Limites tensorial sem massa, quase nulo e próximo ao limiar &
Recuperar resultados pertinentes, explicitando setores e convenções. \\ \addlinespace[4pt]
Numérico & Integração angular e truncamentos espectrais & Tolerância
inicial de \(10^{-3}\) nas ORFs, com erro absoluto perto de zeros;
ajustar para ficar abaixo do erro estatístico relevante. \\ \addlinespace[4pt]
Inferência & Recuperar massa e amplitudes injetadas & Medir viés
normalizado, largura dos intervalos e estabilidade entre amostradores ou
execuções. \\ \addlinespace[4pt]
Calibração & Injeções com parâmetros sorteados da priori e testes em
pontos fixos & Verificar calibração sob a priori; medir separadamente
cobertura em parâmetros fixos, sem exigir que cobertura bayesiana seja
nominal em todo ponto. \\ \addlinespace[4pt]
Detecção & Realizações de RG com ruído e fundos alternativos & Calibrar
a distribuição do fator de Bayes ou estatística escolhida sob a hipótese
nula. \\ \addlinespace[4pt]
Informação & Comparar priori e posterior & Medir contração, divergência
de informação quando apropriada e dependência dos limites com o suporte
da priori. \\ \addlinespace[4pt]
Robustez & Mudar espectro, ruído, distâncias e covariâncias & Apresentar
quais conclusões resistem e onde a aproximação falha. \\
\end{longtable}
\endgroup
}

Para cobertura, um conjunto de 500 realizações teria incerteza binomial
aproximada de 1,3 ponto percentual em torno de 90\%. Esse número é um
alvo de planejamento para cenários selecionados, não uma promessa de
executar toda a grade com 500 repetições. Taxas de falso positivo da
ordem de 1\% requerem amostras maiores ou intervalos de incerteza
explicitamente largos.

